% GENERATED FILE. Edit canonical lessons, then run npm run generate:books.
% canonical-source-hash: fnv1a64:a5c578689d447a34
% canonical-lessons: ES-C647-semaforo, ES-C647-bocacalle, ES-C647-calzada, ES-C647-peaton, ES-C647-rotonda, ES-R647-repaso-640-647

\chapter{Cruzar la calle}
\label{ch:es-cruzar-la-calle}
\hlchaptermodality{647}

\begin{chapteropening}
\textbf{By the end of this chapter:} \emph{I can get across town safely: traffic lights, turnings, the roadway, pedestrians and roundabouts.}

Complete the last lesson of chapter 647: I can get across town safely: traffic lights, turnings, the roadway, pedestrians and roundabouts.
\end{chapteropening}

\section[el semáforo]{el semáforo — the traffic lights}

\label{lesson:ES-C647-semaforo}



\begin{quote}
Before the new one: say the Spanish for the fuel, then the Spanish for the kilometre.
\end{quote}

\subsection*{You'll want to know: el semáforo}
\textbf{el semáforo} — \textquotedblleft{}the traffic lights\textquotedblright{}.

\begin{quote}
Espera, el semáforo está en rojo.
\end{quote}

\begin{quote}
\emph{Wait, the lights are red.}
\end{quote}

One word for the whole set of lights: \emph{el semáforo}.

\begin{cousinweb}
\emph{semáforo} was borrowed from French \emph{sémaphore} (EtymDB), a word made from Greek \emph{sêma} \textquotedblleft{}sign\textquotedblright{} and \emph{-phóros} \textquotedblleft{}carrying\textquotedblright{}: a thing that carries signals.
\end{cousinweb}

\begin{grammarlens}[title={What you've built}]
One more word for telling, explaining and describing at B1.
\end{grammarlens}

\subsection*{Your turn}
\begin{itemize}
\raggedright
  \item \emph{Say it:} \emph{el semáforo}
  \item \emph{Say it:} \emph{el semáforo}, once more
  \item \emph{Say it:} a sentence of your own with \emph{el semáforo}
  \item \emph{Recall:} say the Spanish for the fuel, then the Spanish for the kilometre, then say \emph{el semáforo} again
\end{itemize}

\begin{tcolorbox}[breakable,colback=teal!4,colframe=teal!35!black,title={Before you move on}]
What does el semáforo mean? (\textquotedblleft{}The traffic lights\textquotedblright{}.) Say it once more.
\end{tcolorbox}
\section[la bocacalle]{la bocacalle — the turning, the side-street opening}

\label{lesson:ES-C647-bocacalle}



\begin{quote}
Before the new one: say the Spanish for the kilometre, then the Spanish for the traffic lights.
\end{quote}

\subsection*{You'll want to know: la bocacalle}
\textbf{la bocacalle} — \textquotedblleft{}the turning, the side-street opening\textquotedblright{}.

\begin{quote}
Gire en la segunda bocacalle a la derecha.
\end{quote}

\begin{quote}
\emph{Take the second turning on the right.}
\end{quote}

Built from \emph{boca} and \emph{calle}: the mouth of a street.

\begin{cousinweb}
\emph{bocacalle} is built from \emph{boca} \textquotedblleft{}mouth\textquotedblright{}, inherited from Latin \emph{bucca} \textquotedblleft{}cheek\textquotedblright{}, and \emph{calle} \textquotedblleft{}street\textquotedblright{}, from Latin \emph{callis} \textquotedblleft{}footpath\textquotedblright{}.
\end{cousinweb}

\begin{grammarlens}[title={What you've built}]
One more word for telling, explaining and describing at B1.
\end{grammarlens}

\subsection*{Your turn}
\begin{itemize}
\raggedright
  \item \emph{Say it:} \emph{la bocacalle}
  \item \emph{Say it:} \emph{la bocacalle}, once more
  \item \emph{Say it:} a sentence of your own with \emph{la bocacalle}
  \item \emph{Recall:} say the Spanish for the kilometre, then the Spanish for the traffic lights, then say \emph{la bocacalle} again
\end{itemize}

\begin{tcolorbox}[breakable,colback=teal!4,colframe=teal!35!black,title={Before you move on}]
What does la bocacalle mean? (\textquotedblleft{}The turning, the side-street opening\textquotedblright{}.) Say it once more.
\end{tcolorbox}
\section[la calzada]{la calzada — the roadway, the carriageway}

\label{lesson:ES-C647-calzada}



\begin{quote}
Before the new one: say the Spanish for the traffic lights, then the Spanish for the turning.
\end{quote}

\subsection*{You'll want to know: la calzada}
\textbf{la calzada} — \textquotedblleft{}the roadway, the carriageway\textquotedblright{}.

\begin{quote}
No juegues en la calzada.
\end{quote}

\begin{quote}
\emph{Do not play in the road.}
\end{quote}

The part of the street for cars; people walk beside it.

\begin{cousinweb}
\emph{calzada} is usually traced to Vulgar Latin \emph{(via) calciāta} \textquotedblleft{}a paved road\textquotedblright{}, built on Latin \emph{calx, calcis} \textquotedblleft{}limestone\textquotedblright{} (Lewis \& Short); the details are not settled.
\end{cousinweb}

\begin{grammarlens}[title={What you've built}]
One more word for telling, explaining and describing at B1.
\end{grammarlens}

\subsection*{Your turn}
\begin{itemize}
\raggedright
  \item \emph{Say it:} \emph{la calzada}
  \item \emph{Say it:} \emph{la calzada}, once more
  \item \emph{Say it:} a sentence of your own with \emph{la calzada}
  \item \emph{Recall:} say the Spanish for the traffic lights, then the Spanish for the turning, then say \emph{la calzada} again
\end{itemize}

\begin{tcolorbox}[breakable,colback=teal!4,colframe=teal!35!black,title={Before you move on}]
What does la calzada mean? (\textquotedblleft{}The roadway, the carriageway\textquotedblright{}.) Say it once more.
\end{tcolorbox}
\section[el peatón]{el peatón — the pedestrian}

\label{lesson:ES-C647-peaton}



\begin{quote}
Before the new one: say the Spanish for the turning, then the Spanish for the roadway.
\end{quote}

\subsection*{You'll want to know: el peatón}
\textbf{el peatón} — \textquotedblleft{}the pedestrian\textquotedblright{}.

\begin{quote}
Los peatones tienen prioridad.
\end{quote}

\begin{quote}
\emph{Pedestrians have priority.}
\end{quote}

\emph{El paso de peatones} — the pedestrian crossing.

\begin{cousinweb}
\emph{peatón} was borrowed from French \emph{piéton} (EtymDB), built on \emph{pied} \textquotedblleft{}foot\textquotedblright{}, which comes from Latin \emph{pēs, pedis}.
\end{cousinweb}

\begin{grammarlens}[title={What you've built}]
One more word for telling, explaining and describing at B1.
\end{grammarlens}

\subsection*{Your turn}
\begin{itemize}
\raggedright
  \item \emph{Say it:} \emph{el peatón}
  \item \emph{Say it:} \emph{el peatón}, once more
  \item \emph{Say it:} a sentence of your own with \emph{el peatón}
  \item \emph{Recall:} say the Spanish for the turning, then the Spanish for the roadway, then say \emph{el peatón} again
\end{itemize}

\begin{tcolorbox}[breakable,colback=teal!4,colframe=teal!35!black,title={Before you move on}]
What does el peatón mean? (\textquotedblleft{}The pedestrian\textquotedblright{}.) Say it once more.
\end{tcolorbox}
\section[la rotonda]{la rotonda — the roundabout}

\label{lesson:ES-C647-rotonda}



\begin{quote}
Before the new one: say the Spanish for the roadway, then the Spanish for the pedestrian.
\end{quote}

\subsection*{You'll want to know: la rotonda}
\textbf{la rotonda} — \textquotedblleft{}the roundabout\textquotedblright{}.

\begin{quote}
En la rotonda, toma la segunda salida.
\end{quote}

\begin{quote}
\emph{At the roundabout, take the second exit.}
\end{quote}

In Spain also \emph{la glorieta}.

\begin{cousinweb}
\emph{rotonda} was borrowed from Italian \emph{rotonda} \textquotedblleft{}a round building\textquotedblright{}, from Latin \emph{rotundus} \textquotedblleft{}round\textquotedblright{} (Lewis \& Short), which is built on \emph{rota} \textquotedblleft{}wheel\textquotedblright{}.
\end{cousinweb}

\begin{grammarlens}[title={What you've built}]
One more word for telling, explaining and describing at B1.
\end{grammarlens}

\subsection*{Your turn}
\begin{itemize}
\raggedright
  \item \emph{Say it:} \emph{la rotonda}
  \item \emph{Say it:} \emph{la rotonda}, once more
  \item \emph{Say it:} a sentence of your own with \emph{la rotonda}
  \item \emph{Recall:} say the Spanish for the roadway, then the Spanish for the pedestrian, then say \emph{la rotonda} again
\end{itemize}

\begin{tcolorbox}[breakable,colback=teal!4,colframe=teal!35!black,title={Before you move on}]
What does la rotonda mean? (\textquotedblleft{}The roundabout\textquotedblright{}.) Say it once more.
\end{tcolorbox}
\section[(dialogue)]{Viajes y ciudades}

\label{lesson:ES-R647-repaso-640-647}



\begin{quote}
Say the Spanish for to spend the summer and the car.
\end{quote}

\subsection*{Your turn}
Say these aloud:

\begin{itemize}
\raggedright
  \item to spend the summer, the car, the traffic, the van, the driver
  \item the taxi, the taxi driver, the fare, the journey, to stop
  \item the ring road, the hard shoulder, the toll, the pothole, the stretch
  \item to skid, to start, to get stuck, to speed up, to circulate
  \item the engine, the steering wheel, the boot, the belt, to fasten
  \item the puncture, to break down, the mechanic, the tyre, to get a puncture
  \item the petrol pump, to refuel, the tank, the fuel, the kilometre
  \item the traffic lights, the turning, the roadway, the pedestrian, the roundabout
\end{itemize}

\begin{tcolorbox}[breakable,colback=teal!4,colframe=teal!35!black,title={Before you move on}]
To spend the summer? (\textbf{veranear}.) The car? (\textbf{el automóvil}.) The traffic? (\textbf{la circulación}.) The van? (\textbf{la furgoneta}.) The driver? (\textbf{el conductor}.) The taxi? (\textbf{el taxi}.) The taxi driver? (\textbf{el taxista}.) The fare? (\textbf{la tarifa}.) The journey? (\textbf{el trayecto}.) To stop? (\textbf{parar}.) The ring road? (\textbf{la circunvalación}.) The hard shoulder? (\textbf{el arcén}.) The toll? (\textbf{el peaje}.) The pothole? (\textbf{el bache}.) The stretch? (\textbf{el tramo}.) To skid? (\textbf{derrapar}.) To start? (\textbf{arrancar}.) To get stuck? (\textbf{atascarse}.) To speed up? (\textbf{acelerar}.) To circulate? (\textbf{circular}.) The engine? (\textbf{el motor}.) The steering wheel? (\textbf{el volante}.) The boot? (\textbf{el maletero}.) The belt? (\textbf{el cinturón}.) To fasten? (\textbf{abrocharse}.) The puncture? (\textbf{el pinchazo}.) To break down? (\textbf{averiarse}.) The mechanic? (\textbf{el mecánico}.) The tyre? (\textbf{el neumático}.) To get a puncture? (\textbf{pinchar}.) The petrol pump? (\textbf{el surtidor}.) To refuel? (\textbf{repostar}.) The tank? (\textbf{el depósito}.) The fuel? (\textbf{el combustible}.) The kilometre? (\textbf{el kilómetro}.) The traffic lights? (\textbf{el semáforo}.) The turning? (\textbf{la bocacalle}.) The roadway? (\textbf{la calzada}.) The pedestrian? (\textbf{el peatón}.) The roundabout? (\textbf{la rotonda}.)
\end{tcolorbox}
